\section{Kết quả thực nghiệm (Pha 1)}
\label{sec:experiments}

\subsection{Thiết lập}
\label{sec:setup}

Bộ kịch bản v0 được sinh tất định từ 130 topology Topology Zoo có 15--40 nút, có toạ độ và liên thông. Các topology được chia tầng theo số nút; mười topology được dùng để đánh giá với ba replicate mỗi topology, ba topology khác dành cho phát triển với hai replicate và không được dùng trong Pha 1, vì các phương pháp cơ sở không có tham số cần hiệu chỉnh. Kích thước cảnh báo lấy tỷ lệ tương đối từ demand của instance SNDlib \texttt{abilene}~\cite{orlowski2010sndlib}; vùng hư hại, điểm nguồn, thời điểm phát sinh, hạn và cạnh hỏng được sinh bằng các luồng số ngẫu nhiên có seed riêng. Họ v0-bh50 giữ nguyên mọi seed và chỉ giảm dung lượng backhaul, nên hai họ cùng topology, cảnh báo và mẫu cạnh hỏng. \Cref{tab:params} liệt kê các tham số. Mọi kết hợp là kịch bản tổng hợp từ dữ liệu công khai, không phải mạng cứu hộ đo thực địa.

Mỗi phương pháp lập kế hoạch một lần cho mỗi scenario và được đánh giá trên 30 realization kênh dùng chung số ngẫu nhiên; cận trên được giải cho từng realization. Trên mỗi họ, kết quả được lấy trung bình theo realization rồi theo \ResultScenarioCount{} scenario đánh giá; khoảng tin cậy 95\% dùng bootstrap theo cụm topology, vì các replicate của cùng một topology không độc lập. Kiểm định ghép cặp giữa các phương pháp thuộc Pha 2. Thí nghiệm chạy bằng Python \ResultPythonVersion, \texttt{scipy} \ResultScipyVersion{} và HiGHS \ResultHighsVersion{} với \ResultWorkers{} tiến trình trên máy \ResultCpuCount{} CPU; tổng thời gian tính của các task là \ResultComputeMinutes{} phút.

% Bảng tham số — hàng đọc từ data/params.csv (sinh bởi experiments/report_tables.py).
\begin{table}[htbp]
\centering
\small
\setlength{\tabcolsep}{4pt}
\caption{Tham số của bộ scenario v0, họ v0-bh50 và thí nghiệm Pha 1.}
\label{tab:params}
\begin{tabular}{@{}lll@{}}
\toprule
\textbf{Tham số} & \textbf{Ký hiệu} & \textbf{Giá trị} \\
\midrule
\csvreader[separator=semicolon, late after line=\\]{data/params.csv}{1=\colparameter, 2=\colsymbol, 3=\colvalue}{\colparameter & \colsymbol & \colvalue}
\bottomrule
\end{tabular}
\end{table}


\subsection{Kết quả chính}
\label{sec:results-main}

% Kết quả Pha 1 — hàng đọc từ data/phase1_main.csv (sinh bởi experiments/report_tables.py).
\begin{table}[htbp]
\centering
\small
\setlength{\tabcolsep}{4pt}
\caption{Kết quả Pha 1 trên 30 scenario đánh giá, mỗi scenario 30 realization kênh. Khoảng tin cậy 95\% lấy bằng bootstrap theo topology; gap tính so với cận trên cùng thông tin (B0: cận chỉ mặt đất, B1: cận trên quỹ đạo B1); thời gian tính bằng giây.}
\label{tab:phase1-main}
\begin{tabular}{@{}llccccccc@{}}
\toprule
\textbf{Backhaul} & \textbf{PP} & \textbf{Đúng hạn} & \textbf{KTC 95\%} & \textbf{Conn} & \textbf{Cận trên} & \textbf{Gap} & \textbf{Lập KH} & \textbf{LP} \\
\midrule
\csvreader[separator=semicolon, late after line=\\]{data/phase1_main.csv}{2=\colbackhaul, 3=\colmethod, 4=\coltimely, 5=\colci, 6=\colconn, 7=\colbound, 8=\colgap, 9=\colplanning, 11=\colboundtime}{\colbackhaul & \colmethod & \coltimely & \colci & \colconn & \colbound & \colgap & \colplanning & \colboundtime}
\bottomrule
\end{tabular}
\end{table}


\Cref{tab:phase1-main} tóm tắt kết quả. Trên họ v0, B1 giao đúng hạn trung bình \ResultTimelyVzeroBone{} số cảnh báo, so với \ResultTimelyVzeroBzero{} của B0; B1 cao hơn B0 trên \ResultBoneWinsVzero{} trong \ResultScenarioCount{} scenario và bằng B0 trên \ResultBoneTiesVzero{} scenario. Khả năng kết nối cho thấy vai trò của UAV rõ hơn: không có UAV, chỉ \ResultConnVzeroBzero{} nguồn có đường theo thời gian tới trung tâm, còn với quỹ đạo B1 tỷ lệ này là \ResultConnVzeroBone. Khoảng tin cậy của hai phương pháp chồng lấn; kết luận thống kê về chênh lệch cần kiểm định ghép cặp ở Pha 2.

Gap của B0 so với cận chỉ dùng mặt đất là \ResultGapVzeroBzero, còn gap của B1 so với cận trên quỹ đạo B1 là \ResultGapVzeroBone. Hai cận dùng thông tin khác nhau, nên gap đo phần còn có thể cải thiện trong cùng điều kiện của từng phương pháp, không dùng để so trực tiếp hai phương pháp. Với quỹ đạo B1 cố định, cận trên cho thấy việc chọn đường và cấp băng thông tốt hơn có thể giao tới \ResultBoundVzeroBone{} số cảnh báo.

Khi backhaul chỉ còn \qty{50}{\kilo\bit\per\second}, tỷ lệ đúng hạn giảm còn \ResultTimelyVzeroBhfivezeroBzero{} với B0 và \ResultTimelyVzeroBhfivezeroBone{} với B1, trong khi gap tăng lên \ResultGapVzeroBhfivezeroBzero{} và \ResultGapVzeroBhfivezeroBone. B1 vẫn cao hơn B0 trên \ResultBoneWinsVzeroBhfivezero{} scenario. Gap tăng cho thấy quy tắc chia đều băng thông và chọn đường theo độ trễ ước tính của các phương pháp cơ sở không thích nghi với nghẽn backhaul, là tình huống mà cơ chế sửa lịch ở Pha 2 nhắm tới.

\subsection{Độ chặt của cận trên}
\label{sec:results-milp}

% Đối chiếu MILP — hàng đọc từ data/phase1_milp.csv (sinh bởi experiments/report_tables.py).
\begin{table}[htbp]
\centering
\small
\setlength{\tabcolsep}{4pt}
\caption{Đối chiếu trên instance thu nhỏ (10 cảnh báo, quỹ đạo B1, realization 0): số cảnh báo đúng hạn của B1 và của kế hoạch dựng từ nghiệm MILP (giá trị tốt hơn giữa băng thông của nghiệm và chia đều theo backlog), nghiệm MILP tốt nhất, cận đối ngẫu MILP, cận LP và thời gian giải MILP.}
\label{tab:phase1-milp}
\begin{tabular}{@{}lccccccc@{}}
\toprule
\textbf{Scenario} & \textbf{$|A|$} & \textbf{B1} & \textbf{KH MILP} & \textbf{MILP} & \textbf{Cận MILP} & \textbf{Cận LP} & \textbf{Thời gian (s)} \\
\midrule
\csvreader[separator=semicolon, late after line=\\]{data/phase1_milp.csv}{1=\colscenario, 2=\colalerts, 3=\colmethod, 4=\colplan, 5=\colincumbent, 6=\coldual, 7=\collp, 9=\colruntime}{\texttt{\colscenario} & \colalerts & \colmethod & \colplan & \colincumbent & \coldual & \collp & \colruntime}
\bottomrule
\end{tabular}
\end{table}


\Cref{tab:phase1-milp} đối chiếu cận LP với MILP trên \ResultMilpCount{} instance thu nhỏ. HiGHS giải tối ưu \ResultMilpOptimalCount{} trong \ResultMilpCount{} bài MILP. Trên các instance này, B1 đạt đúng giá trị tối ưu MILP ở \ResultMilpMethodAtOptimumCount{} instance và kế hoạch dựng từ nghiệm MILP đạt tối ưu ở \ResultMilpPlanAtOptimumCount{} instance, nên khoảng cách giữa cận LP và giá trị đạt được chủ yếu đến từ việc nới lỏng tính nguyên của lựa chọn đường. Khi dựng kế hoạch, băng thông lấy thẳng từ nghiệm MILP chỉ giao đúng hạn tổng cộng \ResultMilpStaticTotal{} cảnh báo trên năm instance, so với \ResultMilpEqualSplitTotal{} khi giữ các đường của nghiệm và chia đều băng thông theo backlog, vì lịch băng thông tĩnh lệch nhịp với cách bộ điều phối EDF phục vụ hàng đợi. Bao mười tiếp tuyến ước lượng tốc độ cao hơn thực tế tối đa \ResultTangentOverestimateMax{} (theo tỷ lệ tương đối) trên $[B_{\mathrm{tot}}/512, B_{\mathrm{tot}}]$. Vì vậy gap trên các scenario đầy đủ ở \cref{tab:phase1-main} nên được hiểu là cận trên của phần có thể cải thiện, có thể lớn hơn đáng kể so với phần thực sự đạt được.

\subsection{Thời gian tính toán}
\label{sec:results-runtime}

Lập kế hoạch B0 gần như tức thời, còn B1 mất trung bình \ResultPlanningVzeroBone{} giây cho mỗi scenario v0 vì phải mô phỏng từng candidate; đánh giá một kế hoạch trên một realization mất khoảng \ResultEvaluationVzeroBone{} giây. \Cref{fig:lp-runtime} cho thấy thời gian giải của \ResultLpCount{} bài LP: số biến có trung vị \ResultLpVariablesMedian{} và lớn nhất \ResultLpVariablesMax{}, thời gian giải có trung vị \ResultLpRuntimeMedian{} giây và lớn nhất \ResultLpRuntimeMax{} giây. Cận trên quỹ đạo B1 trong họ v0-bh50 mất trung bình \ResultBoundTimeVzeroBhfivezeroBone{} giây, lâu hơn \ResultBoundTimeVzeroBone{} giây của họ v0.

\begin{figure}[htbp]
\centering
\begin{tikzpicture}
\begin{loglogaxis}[
    width=0.8\textwidth, height=6cm,
    xlabel={Số biến của bài toán LP}, ylabel={Thời gian giải (s)},
    grid=major, /pgf/number format/use comma,
]
\addplot[only marks, mark=*, mark size=0.5pt, color=blue!60!black, opacity=0.4]
    table[col sep=comma, x=variables, y=runtime] {data/lp_runtime.csv};
\end{loglogaxis}
\end{tikzpicture}
\caption{Thời gian giải của \ResultLpCount{} bài toán LP cận trên (hai họ scenario, hai loại cận, 30 realization) theo số biến.}
\label{fig:lp-runtime}
\end{figure}
